%% Chapter 12: Approximation Algorithms I
%% Status: skeleton -- learning goals and section outline fixed, prose to be written.

\chapter{Approximation Algorithms I}
\label{ch:approx-i}

An approximation algorithm runs in polynomial time and returns a solution
together with a proof that it is not far from optimal. The art is in the proof:
we must compare against an optimum we cannot compute, so every approximation
guarantee rests on a lower bound (or upper bound) on $\OPT$ that we can reason
about.

\section*{Learning goals}
\addcontentsline{toc}{section}{Learning goals}

After this chapter you should be able to:
\begin{itemize}[leftmargin=*]
  \item define approximation ratios for minimisation and maximisation problems;
  \item analyse greedy approximation algorithms and prove their ratios;
  \item give and analyse the $2$-approximation and Christofides' algorithm for metric \lang{TSP};
  \item adapt the metric \lang{TSP} arguments to a near-symmetric asymmetric instance.
\end{itemize}

\section{What an approximation guarantee is}
\label{sec:approx-definitions}

\todo{Write this section.}
% Planned content: Absolute versus relative error, and the role of bounds on $\OPT$.

\section{Greedy algorithms}
\label{sec:greedy-approx}

\todo{Write this section.}
% Planned content: Vertex cover in factor $2$; set cover in factor $\ln n$.

\section{Local search}
\label{sec:local-search}

\todo{Write this section.}

\section{Metric TSP}
\label{sec:metric-tsp}

\todo{Write this section.}
% Planned content: The double-tree $2$-approximation and Christofides' $3/2$-approximation.

\section{Asymmetric instances}
\label{sec:asymmetric-tsp}

\todo{Write this section.}
% Planned content: When $d(i,j) \le 2 d(j,i)$ the metric arguments survive at a constant loss.

\section{Scheduling and load balancing}
\label{sec:scheduling-approx}

\todo{Write this section.}

\section{Exercises}
\label{sec:ex-approx-i}

\todo{Write this section.}

\begin{poolquestions}
  \poolq{13}{Approximation non-symmetric TSP}Covered in \cref{sec:asymmetric-tsp}.
\end{poolquestions}
